% ============================================================
%  7. fejezet, A szoftverrendszer rövid ismertetése
% ============================================================
\chapter{A szoftverrendszer rövid ismertetése}
\label{ch:user_docs}

\section{Rendszer áttekintése}

A KCMamba-kutatás kísérleti infrastruktúrájaként egy teljes körű,
produkciókész kereskedési rendszer készült el. A rendszer négy fő részből
áll: (i)~\textbf{aszinkron kereskedési motor} hexagonális architektúrában,
(ii)~\textbf{visszatesztelő keretrendszer} historikus adatokon,
(iii)~\textbf{Binance-adapter} élő piaci adathoz és megbízás-végrehajtáshoz,
valamint (iv)~\textbf{webes irányítópult} a valós idejű megfigyeléshez.

A motor eseményvezérelt, portok és adapterek köré épített architektúrát
(\emph{hexagonal / ports-and-adapters}) alkalmaz~\cite{cockburn2005hexagonal},
amely biztosítja, hogy a \texttt{core} üzleti logika tesztelő és éles
adapterrel egyaránt, változtatás nélkül futtatható. A KCMamba modell
cserélhető: a \texttt{ModelStore} interfészen keresztül bármely
\texttt{nn.Module} betölthető súlyok nélkül vagy előtanított
\texttt{.pt} fájlból.

\begin{figure}[H]
  \centering
  \small
  \begin{tabular}{|p{3.8cm}|p{10.5cm}|}
    \hline
    \textbf{Komponens} & \textbf{Feladat} \\
    \hline\hline
    \texttt{TradingEngine}    & Fő koordinátor; eseményhurkot (asyncio) futtat, állapotot karbantart \\
    \texttt{BinanceFeed}      & WebSocket-csatornán gyertyaadatot és könyvmélységet fogad \\
    \texttt{KCMamba} (SSM)    & 96~lépéses kontextusból 96~lépéses előrejelzést ad \\
    \texttt{ThresholdStrategy}& Előrejelzésből \textsc{Buy/Sell/Hold} jelzést generál \\
    \texttt{RiskPolicy}       & Pozícióméret-korlát, stop-loss, napi veszteség-limit \\
    \texttt{BinanceBroker}    & Futures REST API; paper- és élő módban cserélhető \\
    \texttt{Dashboard}        & aiohttp szerver, 8~nézet, valós idejű SSE frissítés \\
    \hline
  \end{tabular}
  \captionof{table}{A rendszer főbb komponensei és feladataik.}
  \label{tab:components}
\end{figure}

\section{Rendszerkövetelmények és telepítés}

A rendszer Python~3.10+ környezetet igényel; GPU opcionális (CPU-n is fut).
Telepítés:

\begin{verbatim}
python -m venv .venv && source .venv/bin/activate   # Windows: .venv\Scripts\Activate.ps1
pip install -r requirements.txt
cp .env.example .env   # BINANCE_API_KEY, BINANCE_SECRET kitöltése szükséges élő módhoz
\end{verbatim}

\section{Futtatás}

\subsubsection*{Paper-trading demó (alapértelmezett)}

\begin{verbatim}
python main.py
\end{verbatim}


A rendszer Binance WebSocket-adatot fog, KCMamba modellel előrejelez,
ThresholdStrategy jelzéseket generál, de valódi megbízást \emph{nem}
küld.

\subsubsection*{Visszatesztelés}

\begin{verbatim}
python -m adapters.testing.backtest --symbol BTCUSDT --start 2024-01-01 --end 2024-06-01
\end{verbatim}

Binance-kapcsolat nélkül fut; eredményt konzolra és CSV-be ír.

\subsubsection*{Offline benchmark (kísérleti kiértékelés)}

\begin{verbatim}
python tdk_sweep.py --resume    # TDK sweep: Exchange, ETTh1, ETTh2 + szintetikus zajok
\end{verbatim}

Ez a parancs reprodukálja a~\ref{ch:kiserletek}.~fejezetben bemutatott
összes kísérleti eredményt; a kimenet \texttt{tdk\_results.json}-ba kerül.

\section{Biztonsági megfontolások}

A demó alapértelmezés szerint paper-trading módban fut, valódi megbízást
nem küld. Éles kulcs esetén a Binance-oldalon (a)~csak olvasási és
Futures-kereskedési jogot adjunk, és (b)~IP-whitelistet alkalmazzunk.
Az irányítópult alapértelmezetten \texttt{localhost:8080}-on hallgat;
nyilvános IP-n csak fordított proxy (nginx) mögött ajánlott.
